% Guide paragraphs are added when the relevant section has been visually collated.
% Cross-references to later sections will be extended as their review progresses.
\paragraph{\S67. Aspheric maps replace contractibility as the organizing notion (pp.\ 188--191).}
Grothendieck reorients the modelizer program: the crucial objects are aspheric objects and weak equivalences, and its natural comparison maps are aspheric functors. He explicitly judges the earlier restriction to contractibility and special test functors to have been an unnecessarily long detour (p.\ 188), while retaining its possible relevance for ``canonical modelizers.'' The discussion now allows arbitrary small categories, rather than only test categories. For a basic localizer $\underline W$ he introduces $(\mathrm{Hot}^{W}_{A})$, also written $(\mathrm{Hot}_{A})$, as the localization of the presheaf category by $\underline W_A$; the construction is then organized by inverse image $i^*$ and the comparison morphism $\lambda_i(F):A_{/i^*F}\longrightarrow B_{/F}$, interpreted through the corresponding induced topoi (pp.\ 189--190). This prepares the promised description of asphericity structures and testing functors at the end of the section (p.\ 191); the next section interrupts that development. The dated marginal correction of 3 May on p.\ 189 expressly retracts the claim that Quillen's weak equivalences are precisely $\underline W_A$, saying that Quillen's are stronger. It is therefore essential to read the original assertion together with the margin. The witness changes handwritten ``or'' into equality in (1), drops the subscript $A$ on $\underline W_A$, omits the handwritten domain and codomain after the action on $(a,p)$ on p.\ 190, and prints the struck word ``real'' on p.\ 188. These differences concern notation, functorial context, or the author's revision; none licenses an unrecorded mathematical repair of the typescript.

\paragraph{\S68. Mathematical work, the Teichmüller tower, and the necessity of groupoids (pp.\ 191--195).}
Dated 27 May, the section explains the six-week interruption in the asphericity discussion through a retrospective of the author's changing mathematical interests. Its mathematical center is the ``Lego-Teichmüller construction game'': cutting and gluing surfaces, the tower of groupoids $T_{g,\nu}$, and the arithmetic action on their profinite completions (pp.\ 191--192, 194--195). The account of \emph{La Longue Marche}, tame topology and ``unfolding'' of stratified structures supplies the context for the anabelian project; the intervening reflection on creative interest, ambition and alternating periods of work belongs to the argument's historical presentation and must not be omitted as an aside (p.\ 193). The final page gives the technical lesson: reducing several natural base points to one can destroy visible symmetries and obscure a presentation by generators and relations, whereas groupoids preserve the geometry. Van Kampen situations supply the comparison, and the modular ``multiplicities'' are identified as topoi whose points have finite automorphism groups (p.\ 195). This is a return to the text's groupoid motivation, rather than a proof of the asphericity claims left pending in \S67; \S69 resumes a retrospective of the technical scratchwork. The section explicitly revises an earlier preference for groups or ``exterior groups'' in favor of groupoids (p.\ 195). The witness has several ordinary wording errors here---including ``case'' for ``care'' on p.\ 195---while the underlining and handwritten exclamation in ``topoi!'' preserve the emphasis of the typescript.

\paragraph{\S69. Derivators and integration of homotopy types (pp.\ 196--202).}
The technical retrospective explains why an ordinary derived or homotopy category retains too little structure: traces, determinants and the compatibility of iterated constructions motivate the system of localized diagram categories $D_A(I)$, contravariant in the indexing category $I$, with inverse-image functors expected to admit both adjoints (pp.\ 196--198). This develops the modelizer program recalled in \S67 and connects it with the geometric motivations of \S68. For categories, a covariant diagram can be integrated by forming its associated cofibered category; the comparison with localized slice categories yields the model for ``integration of homotopy types'' (pp.\ 198--199). The name ``derivator'' is introduced on p.\ 200, together with the internal stability expressed by $D(I\times J)$ and the aspiration to a corresponding homotopy theory of derivators. The geometric example is the system of tubes around strata of Mumford--Deligne compactifications: their $K(\pi,1)$ homotopy types, with Teichmüller-type groups, suggest an algebraic description from a diagram of fundamental groupoids (p.\ 201). A diagram in $(\mathrm{Hot})$ alone is explicitly said to be insufficient; a diagram in a model category is required. Grothendieck retracts his recent expectation that finite limits exist in $(\mathrm{Hot})$ (p.\ 201), and \S70 immediately examines the extra factorization condition proposed for the two adjoints (pp.\ 200, 202). Significant witness differences include $T$ replaced by $\mathscr D$ in the statement about inverse limits (p.\ 198), $q=f\circ p$ replaced by an equation beginning with $g$, and ``basic modelizer'' replaced by ``basic localizer'' (p.\ 199). These readings are clear on the images and are preserved rather than repaired from the witness.

\paragraph{\S70. Base change, duality, and the limitations of $(\mathrm{Cat})$ (pp.\ 202--206).}
Starting from the desired two-stage factorization $X\to Z\to Y$, Grothendieck asks for preservation of weak equivalences by base change along the second map and co-base change along the first, with either factor optionally a weak equivalence (p.\ 202). The goal supplies the extra condition proposed for derivator operations in \S69. The investigation of the base-change side leads to cohomological smoothness and properness of functors, characterized through asphericity; fibered categories are smooth, cofibered categories proper, and passage to opposite categories interchanges the two conditions (p.\ 203). The pairing of presheaves and copresheaves suggests a broader duality of topoi, but the claim that cosheaves on an arbitrary topos form a topos remains a question (pp.\ 203--204). The co-base-change side is less tractable: open or closed immersions in $(\mathrm{Cat})$ fail the hoped-for general assertion, whereas pushouts along monomorphisms in $A^{\wedge}$ have the required homotopy meaning through the integral of the associated span (pp.\ 204--205). The stated smooth/proper exceptions and the possible special case of an immersion that is also a weak equivalence must be distinguished from that rejected general assertion. This forces a return to elementary modelizers and Quillen factorization questions; the author finally accepts the Kan condition after describing his resistance to it, while retaining dissatisfaction with the reduction to semisimplicial methods (p.\ 206). These questions become part B of the program in \S71. The witness omits the monomorphism hook in the p.\ 204 square and replaces the actual objects $Y,X,Y'$ of the p.\ 205 integral diagram by categories of elements. The scan-confirmed diagrams preserve the distinction between these two levels of the construction.

\paragraph{\S71. A provisional program: modelizers, derivators, and topoi (pp.\ 206--212).}
The program dated 30 May organizes the preceding scratchwork into four tasks. Part A returns to the asphericity structures and canonical modelizers suspended in \S67; part B develops the base-change, co-base-change and factorization questions of \S70, with homotopy integrals, amalgamated sums and Mayer--Vietoris sequences as the connecting constructions (pp.\ 206--207). A handwritten transposition on p.\ 207 makes strong saturation one of the properties expected to follow from the first three basic-localizer axioms. Part C develops the derivator proposal of \S69 while asking what additional structure is needed for external Hom's and basic invariants. It then distinguishes group objects modeling pointed connected homotopy types from the truncated output of group objects in $(\mathrm{Cat})$, and asks which modelizers satisfy a Dold--Puppe theorem (pp.\ 207--209). The proposed trilogy consists of semisimplicial, cubical and hemispherical complexes, the last also called globular complexes by Brown; their boundary and degeneracy operations are related to regular cellular subdivisions of spheres. Multicomplexes give multiple chain complexes, so the simple Dold--Puppe statement does not extend unchanged. The classification of the three series and the need for enough degeneracy maps are presented as questions to investigate, not proved theorems (pp.\ 209--210). Part D, ``Back to topoi,'' seeks relative homotopy categories $D(T)$ and their possible opposite counterparts, extending the $f^*,f_!,f_*$ formalism from indexed categories to topoi. The discussion keeps existence of derived adjoints tentative and moves from objects to proobjects when a topos fails to determine an essentially constant homotopy type (pp.\ 210--211). Local asphericity then reconnects the program with \S35; the revised distinction between modelizing and strictly modelizing topoi is recalled from p.\ 68, and the section ends with the Lawvere-element condition for a finitely triangulated contractible space (p.\ 212). The next section begins the promised return to the asphericity game. The witness changes $f_*$ to $f^*$ and $Lf_!$ to a subscript-$1$ expression on p.\ 210, and applies a handwritten transposition there without marking it. The duplicated target-only display on p.\ 211 was an error of the draft transcription, absent from the scan; its removal restores the actual three-display sequence.

\paragraph{\S72. Criteria for aspheric functors and the first stability results (pp.\ 213--219).}
Returning to \S67 after the program of \S71, Grothendieck isolates the localizer assumptions needed for asphericity: mild saturation, the condition on categories with final object, and the locality condition Loc~3 (pp.\ 213--214). Proposition~1 identifies the comparison maps on categories of elements, their restrictions over objects, and the corresponding slice criteria; it then distinguishes these equivalent asphericity conditions from preservation of weak equivalences and equivalence after localization (pp.\ 214--216). The extra hypotheses in the implication diagram matter: the discrete-target example separates aspheric fibers from merely nonempty fibers, and the product-projection example shows that inducing the localized equivalence need not determine an aspheric map without the stated assumptions (p.\ 216). Under asphericity, the localized triangle commutes up to the canonical comparison isomorphism; this reconnects the discussion with the earlier ``inspiring assumption'' of \S28 without using that assumption as a proof (pp.\ 216--217). Proposition~2 transfers total asphericity along an aspheric functor and relates the inverse image of a separating homotopy interval to asphericity of the functor. Its handwritten addition makes the source category aspheric, and the long interlinear insertion specifies the homotopy condition (TH1) of p.\ 50; the corollary deduces strict test categories from a totally aspheric source and a local test target (pp.\ 217--218). Proposition~3 supplies composition and right cancellation, base change along a fibering map, and products (pp.\ 218--219). The base-change proof uses the fibration condition L5, whose derivation from the stated localizer assumptions is expressly postponed; Grothendieck explains that only part~(a) is needed for the immediate asphericity-structure formalism, preventing the announced application from relying on that deferred argument (p.\ 219). Remark~4 denies the opposite cancellation assertion and leads directly to the examples in \S73. The visual collation restores the proof-diagram arrow from (i'') to (iii), rather than to (i'), on p.\ 215; it also identifies the apparent deletion of the t in Hot on p.\ 214 as pen retracing. On p.\ 218 the comparison sign is a handwritten tilde over a bar, with no arrowhead, contrary to the draft transcription and its diff description. These are distinctions between what the image shows and what a mathematically plausible reconstruction might supply.

\paragraph{\S73. Adjoints, coasphericity, and cofinality (pp.\ 220--221).}
The continuation dated 11 June begins by suppressing the qualifying $\underline W$ when one basic localizer is understood. It sharpens Remark~4 of \S72: a functor from the one-point category into an aspheric category is aspheric exactly when its chosen object is initial, illustrating how much stronger a map condition is than asphericity of its source and target (p.\ 220). Proposition~4 enlarges the preceding isomorphism example: every equivalence is aspheric, and more generally a left adjoint is aspheric because its slice categories identify with slices having final objects. Passage to opposite categories then motivates the name ``coaspheric'' for the dual condition, satisfied by the right adjoint. The corresponding undercategories and their initial objects connect coasphericity with a promised all-dimensional cofinality statement, whose dimension-zero consequence is invariance of direct limits (pp.\ 220--221). The analogy from \S70 returns: asphericity resembles properness with aspheric fibers, and coasphericity resembles smoothness with aspheric fibers. The comparison is explicitly weakened because arbitrary base change does not preserve either condition; the applicable base changes are fibering/smooth or cofibering/proper, respectively. Even an equivalence of categories can have an empty fiber, so its simultaneous asphericity and coasphericity do not supply arbitrary base-change stability. Universal weak equivalences, called ``trivial Serre fibrations'' here, are proposed as a stronger common condition and assigned to a later systematic investigation (p.\ 221). The section then turns to the overdue product discussion of \S74. On p.\ 220 the crossed handwritten heading above Indeed is Remark~5, not the draft's tentative Proof; it is written again at the subsequent paragraph. Whether the placement strokes also delete its typed initial The remains unresolved and is marked as such. The p.\ 221 witness changes singular condition to plural conditions and supplies a title for \S74 absent from the typescript.

\paragraph{\S74. Products, diagonal functors, and canonical comparison of models (pp.\ 221--230).}
The overdue product discussion supplies concrete constructions for the abstract stability results of \S72. The external product of presheaves has a category of elements identified with the product of the two categories of elements; it preserves aspheric objects and weak equivalences in the stated variables, and yields the product assertions for locally aspheric and totally aspheric categories (pp.\ 221--222). The corollaries distinguish local test, test, and strict test categories: the hypothesis on the other factor changes from arbitrary to aspheric to totally aspheric. Mixed multicomplexes illustrate the construction (p.\ 223), reconnecting it with the trilogy in \S71. The comparison functor associated with an aspheric map between modelizing categories leads to a functor from $\mathrm{Asph}(A,B)$ into a category of comparison functors and factors through its fundamental groupoid. Nonemptiness and connectedness suffice for the existence of an isomorphism class of comparison, while a canonical transitive system requires the stronger simple-connectedness conclusion; the text carefully separates these claims (pp.\ 224--225). The diagrams on p.\ 226 distinguish categories, presheaf modelizers, and localized homotopy categories, including the middle presheaf triangle omitted from the witness.
The diagonal criterion then identifies total asphericity of $A$ with asphericity of $A\to A\times A$, and permits simultaneous comparison through products of strict test categories (p.\ 227). Proposition~4 treats pointwise products of functors into a category $B$ with binary products and aspheric presheaves $\underline{\mathrm{Hom}}(b,y)$: the total-asphericity assumption on $A$, moved by a handwritten transposition into the hypotheses, is essential to the stated construction (p.\ 228). The external-product corollary is sharpened by revision to an equivalence: when $i$ is aspheric, $i\boxtimes i'$ is aspheric if and only if $A'$ is aspheric. Its proof uses a lemma on a presheaf over $A\times A'$ whose restrictions to $A$ are aspheric. The resulting category-of-elements projection is a weak equivalence by the fibration condition L5, explicitly recalling the claim and deferred justification on p.\ 219; this dependency should remain visible (p.\ 229). Contractibility of the objects of $B$ is discussed as a stronger practical hypothesis. Finally, closure of $\mathrm{Asph}(A,B)$ under binary products makes its nonempty case connected with the needed homotopy strength, giving the canonical transitive system of isomorphisms between the localized functors and the canonical functor announced on p.\ 225 (p.\ 230). A handwritten margin points forward to the generalization in \S82, p.\ 272. The image-based edition retains the unusual $X\longrightarrow B_{/X}$ display on p.\ 224, the missing slash in the final category-of-elements notation on p.\ 226, and the visibly typed product x and lowercase subscript o on p.\ 228. The hemisphere-like symbol on pp.\ 223 and 229 remains a doubtful approximate glyph; neither mathematical expectation nor the witness settles its exact drawn form.

\paragraph{\S75. Asphericity structures, testing functors, and weak equivalences (pp.\ 230--237).}
The generalities of \S\S72--74 now become a structure on a possibly large $\underline U$-category $M$. After making the universe convention explicit, Grothendieck defines a class $M_{as}$ of aspheric objects through a small category $A$ and a functor $i:A\to M$ that lands in $M_{as}$ and detects it by the presheaf $i^*(x)=\mathrm{Hom}(i(-),x)$ (pp.\ 230--231). Such an $i$ is a “testing functor.” Every small full subcategory supplies an asphericity structure, and Proposition~1 plus its corollary shows that any structure can be generated this way; enlarging the generating full subcategory within $M_{as}$ preserves the structure (p.\ 232). The following terminology discussion distinguishes testing functors from the earlier test functors of pp.\ 175--176. The essential condition depends on the possibly large category $M_{as}$ itself, rather than on its embedding in $M$, so the author tentatively adopts $M_{as}$-aspheric, or simply aspheric, as a synonym. His use of both “aspherical” and “aspheric” in this discussion is preserved (p.\ 233).
Proposition~2 defines $W=(i^*)^{-1}(\underline W_A)$ and proves independence of the testing functor by passing to a common small full subcategory. Mild saturation is immediate; strong saturation remains conditional on the corresponding result for the basic localizer (p.\ 234). When $M$ has a final object, calling the structure aspheric means that this final object belongs to $M_{as}$; the resulting criterion recovers aspheric objects from the maps to the final object. The small-generator descriptions then characterize both the admissible weak-equivalence sets and the asphericity classes. Without the aspheric final-object hypothesis, a weak-equivalence set need not determine a unique asphericity structure. The analogy with contractibility and homotopism structures in \S\S51--52 motivates a promised construction from contractibility to asphericity and a uniqueness assertion for the intended canonical examples (p.\ 235). Proposition~4 recasts the criteria of p.\ 214 as six conditions: three about aspheric objects and three about weak equivalences. The object conditions are equivalent and imply the others; the revised converse requires both $A$ and the asphericity structure to be aspheric. All six conditions become equivalent when the specified generating category $B$ contains the final object. If no $B$ is fixed, enlarging one to contain $i(A)$ and the final object yields equivalence of the first five conditions (pp.\ 236--237). These added hypotheses, some typed interlinearly and some handwritten, are central to the proposition and were partly concealed by malformed insertion macros in the draft. Section~76 immediately tests the formalism on examples.
Two visible mathematical slips are retained for the reader to assess: the universe paragraph says “two objects of $U$” where the witness substitutes $M$ (p.\ 230), and Proposition~2 says “Then $M$ is” a mildly saturated subset of $\mathrm{Fl}(M)$ after defining $W$, where the witness prints $W$ (p.\ 234). Proposition~3 and its corollary also call the relevant map “aspheric,” while the explicit revisions of Proposition~4 replace that word by “a weak equivalence” (pp.\ 234, 236); the edition does not silently harmonize the earlier wording. Several correction layers remain uncertain, including the usual-localizer subscript at the preceding section boundary and overtyped letters in these definitions and proofs.

\paragraph{\S76. Basic examples and total asphericity (pp.\ 237--239).}
Three examples test the definitions of \S75. For $M=(\mathrm{Cat})$, the $\underline W$-aspheric categories form the basic asphericity structure, generated through $a\mapsto A_{/a}$ for a weak test category $A$. The shift from the basic modelizer to this structure puts aspheric objects before weak-equivalence arrows. The distinction matters for functors into $(\mathrm{Cat})$: asphericity supplies the induced localized functor, whereas the weak-test requirement asks for a homotopy-model comparison. Theorem~1 of p.\ 176 suggests equivalence of these conditions when the values are contractible, and the author identifies a comprehensive answer to the “silly question” of \S46, p.\ 95, as the motivation for the present formalism (pp.\ 237--238). The second example takes topological spaces with the weakly contractible spaces as aspheric objects; the geometric realization of simplices supplies the testing functor, and the text appeals to Gabriel--Zisman for the stronger modelizer statement. Grothendieck adds the previously omitted assertion that the first two structures are aspheric and later notes that the resulting weak equivalences are the usual ones (p.\ 238). The third example is $A^\wedge$ with its usual aspheric presheaves, generated by the canonical fully faithful inclusion of any small $A$. Its weak equivalences are $\underline W_A$, while its asphericity is equivalent to asphericity of $A$; the modelizing property is not asserted for arbitrary $A$. A functor $A'\to A$ is aspheric precisely when its composite into this presheaf structure is $M_{as}$-aspheric (pp.\ 238--239). This leads to “totally aspheric” structures: $M$ admits finite products and finite products of aspheric objects, including the empty product, are aspheric. Proposition~5 identifies this condition with the existence of a small generating subcategory stable under finite products, and with a testing functor from a category whose presheaf category is totally aspheric. It reconnects the product arguments of \S74 with the general formalism and prepares the localized canonical functor of \S77 (p.\ 239). The visual check restores a missing typed a before the repaired l in “realization,” retains the literal thatnthis and appheric typing, and distinguishes typewritten interlinear additions from pen text. The upward margin arrow at the end of the examples had been misrepresented by three brackets in the draft; its direction is restored, without attributing a mathematical meaning to the layout mark.


\paragraph{\S77. The canonical homotopy functor and the meaning of modelizing (pp.\ 239--241).}
An asphericity structure now determines more than its weak-equivalence class: it supplies a canonical functor from $\mathrm{Hot}_M=W^{-1}M$ to $(\mathrm{Hot})$. Choosing a testing functor $i:A\to M$, Grothendieck uses the composite through $A^\wedge$ and its category-of-elements functor; the inverse-image descriptions of the weak equivalences permit localization, while conservativity is stated conditionally on strong saturation (pp.\ 239--240). To compare two choices, he places their images in one small full subcategory $B$ of $M_{as}$, obtains the comparison isomorphisms from aspheric functors, and then enlarges $B$ to prove independence. A common $B$ for three choices gives transitivity, completing the canonical construction promised in \S74. Formula (13) on p.\ 240 actually names its composite $i_M$, while the following localized functor carries the extra choice index; the witness's harmonization to $i_{M,i}$ in (13) is not adopted. The two sloping comparison signs there are hand-drawn bars with tildes, without arrowheads. The new definition calls the structure modelizing when this particular canonical functor is an equivalence; this is stronger than merely requiring an abstract equivalence between $W^{-1}M$ and $(\mathrm{Hot})$. In the presheaf example it says exactly that $A$ is a pseudo-test category (p.\ 241). The author immediately hesitates over the terminology and proposes a parallel to the hierarchy of pseudo-test, weak test, test, strict test categories and contractors, but leaves the final qualification as question marks. His deletion of “progressively stronger” in that discussion remains visible; the witness prints the deleted words as live text. Section~78 therefore begins with the modelizing name expressly provisional, and uses the canonical construction to revisit the observation and silly question of \S46 rather than treating the proposed hierarchy as already settled.


\paragraph{\S78. Test functors separated into asphericity and modelizing conditions (pp.\ 241--246).}
The provisional definition of modelizing from \S77 supplies the setting for a reformulation of the observation and silly question of \S46. Theorem~1 assumes a modelizing structure $(M,M_{as})$, a pseudo-test category $A$, a functor $i:A\to M$ with aspheric values, and the existence of a final object of $M$; the last hypothesis is a typed interlinear addition, not an original underlined parenthetical sentence. It equates the asphericity and weak-equivalence criteria of Proposition~4, p.\ 236, with the condition that $i^*$ induce an equivalence between the localized categories (p.\ 242). A lemma says that a final object in a modelizing structure is aspheric. Its supporting corollary characterizes $M_{as}$ as the objects whose images under the canonical map to $(\mathrm{Hot})$ are final. The proof uses strong saturation of the basic localizer, announced as a consequence of Loc~1--3 but deferred to part~V; the author explicitly offers the additional asphericity assumptions on the two final objects as a way of avoiding the feared circularity before that proof is supplied (pp.\ 242--243). The canonical triangle from \S77 then proves the localized equivalence. The proof's literal reference to conditions through (ii') is retained even though the statement speaks of six conditions through (ii'').
The theorem prompts a revision of the terminology inherited from \S65. The evolving pseudo-test-functor definition first describes model-preserving composites, then separates the intended conditions into: $A$ is pseudo-test, $(M,M_{as})$ is modelizing, and $i$ is $M_{as}$-aspheric. The discussion revises both its hypotheses and its formulation on the page; these are stages of the argument rather than interchangeable completed definitions (pp.\ 243--244). Passing to weak test functors strengthens the condition on $A$ to weak test. Passing to test functors requires a test category and asphericity both of $i$ and of all its slice-induced functors. Definition~2 calls the latter family of slice conditions locally aspheric, and the conjunction with asphericity of $i$ totally aspheric. The local criterion additionally requires that $i$ factor through $M_{as}$ and that $i^*(x)$ be aspheric over the final presheaf, with $x$ restricted to a generating subcategory containing $i(A)$ when desired (p.\ 245). The Remarks paragraph draws the conceptual conclusion: the independent asphericity conditions have wider relevance than the combined test-functor terminology. Yet the earlier theorem on $(\mathrm{Cat})$ still has one advantage, namely its criterion involving only $i^*(\Delta_1)$ under the contractibility hypothesis on the values of $i$. Recovering such a minimal criterion requires relating asphericity structures to the contractibility structures of \S51, which becomes the subject of \S79 (p.\ 246). The typed minimal is inserted before criterion without deleting generating in the preceding line. Most of the final sentence on p.\ 246 is physically covered by a laid-over sheet; its tentative transcription remains doubtful and is not treated as fully recovered evidence.

